\documentclass[10pt,a4paper]{amsart}
\usepackage[margin=19mm]{geometry}
\usepackage{amsmath,amssymb,mathtools}
\usepackage[scr=rsfs,cal=euler]{mathalfa}
\usepackage{xfrac}
\usepackage[unicode,hidelinks,bookmarks=false]{hyperref}
\newcommand{\ie}{i.e.\ }
\newcommand{\cf}{cf.\ }
\newcommand{\resp}{respectively\ }
\newcommand{\Subsection}{\subsection}
\providecommand{\nobreakdash}{\nobreak\hbox{-}}
\setlength{\emergencystretch}{2em}
\multlinegap0pt
\usepackage{mathtools}
\usepackage{xcolor}
\usepackage[shortlabels]{enumitem}
\usepackage{amssymb}
\usepackage{bm}
\usepackage{tikz}
\usepackage{float}
\newtheorem{theorem}{Theorem}
\newtheorem{proposition}{Proposition}
\def\bz{{\mathbb{Z}}}
\def\gdeg{G\text{\rm -deg}}
\def\wt{\widetilde}
\title{The Ize Conjecture Redux: A Parity Criterion for Global Equivariant Bifurcation Guarantees}
\author{Ziad Ghanem}
\date{Source: arXiv:2606.05425v1 (CC BY 4.0)}
\begin{document}
\maketitle

\noindent\emph{Source and licence.} The Ize Conjecture Redux: A Parity Criterion for Global Equivariant Bifurcation Guarantees. Version arXiv:2606.05425v1 (CC BY 4.0), distributed by arXiv under the Creative Commons Attribution 4.0 International licence, http://creativecommons.org/licenses/by/4.0/, as recorded in the arXiv licence metadata for this exact version and shown on the abstract page https://arxiv.org/abs/2606.05425v1. Full licence notice: \texttt{licenses/arxiv-2606.05425-CC-BY-4.0.txt}.\par\smallskip

\noindent\textit{Editorial scope.} Counted unit: the article's proposition prop:spectral\_flow\_evaluation (source0.tex lines 458--466). The article's complete original proof of it is delivered with it (source0.tex lines 467--499, one \texttt{proof} environment of 33 lines). No mathematics is written, repaired or re-derived: the statement and the proof are the article's own lines, in the article's own notation, and no retained line is substituted or re-flowed.  Retained uncounted as the source-local context the proof needs: lines 351--360; lines 940--940.  Nothing was omitted from any retained span, so the package carries no omission record.  No retained line contains a citation, so the wrapper carries no bibliography and no bibliography entry.  No retained line contains a figure environment, an image-inclusion command or drawing code, so no image is delivered and the package declares no asset dependency.  Editorial formatting only, in addition: an AMS \texttt{amsart} preamble that restates the article's own declarations and macros used by the retained lines, namely the environments \texttt{theorem}, \texttt{proposition} on separate counters, together with the standard packages those lines need.  The retained environments print in this document as Theorem 1, Proposition 1 in order of appearance, whereas the article prints the same items with the numbering of its own full document.  The complete native source of the article as distributed in the arXiv e-print for this exact version is bundled unchanged as \texttt{sources/202114/source0.tex}.
\begin{theorem}[Parity Criterion]\label{thm:parity_guarantee}
An isotropy subgroup $H \leq G$ satisfies the parity condition
\begin{equation*}
    \dim \mathbb V^H - \dim \mathbb V^G \equiv 1 \pmod{2},
\end{equation*}
if and only if there exist an \textbf{odd} number of non-trivial $G$-isotypic indices $i \in \mathcal I \setminus \{0\}$ satisfying
\begin{enumerate}\item[(i)] $m_i \equiv 1 \pmod{2}$ AND
\item[(ii)] $\dim \mathcal V_i^H \equiv 1 \pmod{2}$.
\end{enumerate}
\end{theorem}

\section{The $G$-Equivariant Leray-Schauder Degree} \label{sec:appendix}

\begin{proposition}[Spectral Flow Evaluation]\label{prop:spectral_flow_evaluation}
If $(G, \mathbb V_{[a,b]})$ is an \textbf{Ize pair} with respect to a maximal orbit type $(H) \in \Phi_0(G; \mathbb V_{[a,b]})$, then one has
\[
  \operatorname{coeff}^H \left(
    \gdeg(\mathscr A(a), B(\mathscr H))
    - \gdeg(\mathscr A(b), B(\mathscr H))
  \right) \neq 0.
\]
\end{proposition}

\begin{proof}
Put $\omega([a,b]) := \gdeg(\mathscr A(a), B(\mathscr H)) - \gdeg(\mathscr A(b), B(\mathscr H)) \in A(G)$ and, for each closed subgroup $K \leq G$, denote its local Brouwer degree difference by $\omega^K([a,b]) := \deg(\mathscr A(a)^K, B(\mathscr H^K)) - \deg(\mathscr A(b)^K, B(\mathscr H^K)) \in \bz$. 
By linearity of the coefficient operator, the Recurrence Formula (Appendix~\ref{sec:appendix}) gives
\begin{align*}
\operatorname{coeff}^H ( \omega([a,b]) )  \quad = \frac{\omega^H([a,b]) - \sum_{L > H} \operatorname{coeff}^L ( \omega([a,b]) ) |W(L)| n(H,L)}{|W(H)|}.
\end{align*}
For any subgroup $L \leq G$, let $\mu_L(\alpha) := \sum_{i \in \mathcal I} c_{i}(\alpha) \dim \mathcal V_{i}^L$ denote the dimension of the $L$-fixed-point space of the negative eigenspace at a regular parameter $\alpha$.
Notice that the local degree differences $\omega^K([a,b]) = (-1)^{\mu_K(a)} - (-1)^{\mu_K(b)}$ strictly take values in $\{-2, 0, 2\}$. As such, the recurrence relation operates entirely over even integers. Evaluating the sequence modulo 4 allows us to distinguish between terms that share the same parity of fixed-point dimensions (which vanish) and those with opposite parities (which evaluate to $2$), i.e.
\[
    \omega^K([a,b]) \equiv \begin{cases} 
        0 \pmod 4 & \text{if } \mu_K(a) - \mu_K(b) \equiv 0 \pmod 2, \\
        2 \pmod 4 & \text{if } \mu_K(a) - \mu_K(b) \equiv 1 \pmod 2.
    \end{cases}
\]
Since $(H)$ is a maximal orbit type in $\mathbb V_{[a,b]}$, the $L$-fixed-point space satisfies $\mathbb V_{[a,b]}^L = \mathbb V_{[a,b]}^G$ for all $(L) > (H)$, implying $(\mu_L(a) - \mu_L(b)) \equiv (\mu_G(a) - \mu_G(b)) \pmod{2}$. Under our modulo 4 encoding, this maximality guarantees
\begin{equation}\label{eq:parity_lock_1}
    \omega^L([a,b]) \equiv \omega^G([a,b]) \pmod 4 \quad \text{for all } (L) > (H).
\end{equation}
Conversely, since $(G, \mathbb V_{[a,b]})$ is an Ize pair, the Parity Criterion (Theorem~\ref{thm:parity_guarantee}) implies $\dim \mathbb V_{[a,b]}^H \not\equiv \dim \mathbb V_{[a,b]}^G \pmod 2$. Consequently, $\mu_H(a) - \mu_H(b)$ and $\mu_G(a) - \mu_G(b)$ must have strictly opposite parities, i.e.
\begin{equation}\label{eq:parity_lock_2}
    \omega^H([a,b]) \not\equiv \omega^G([a,b]) \pmod 4.
\end{equation}
We now evaluate the integers $\operatorname{coeff}^L ( \omega([a,b]) ) |W(L)| \pmod 4$ for all $(L) > (H)$ via downward induction. For the top element $G$, the recurrence sum is empty, giving $\operatorname{coeff}^G ( \omega([a,b]) ) |W(G)| = \omega^G([a,b])$. For any intermediate orbit type $H < L < G$, assume inductively that one has $\operatorname{coeff}^{\wt K} ( \omega([a,b]) ) |W(\tilde K)| \equiv 0 \pmod 4$ for all $L < \tilde K < G$. Evaluating the recurrence relation at $L$ modulo 4 yields
\begin{align*}
 \operatorname{coeff}^L ( \omega([a,b]) )|W(L)|  \equiv \omega^L([a,b]) - \sum_{L < \tilde K} \operatorname{coeff}^{\wt K} ( \omega([a,b]) ) |W(\tilde K)| n(L, \tilde K) \pmod 4.    
\end{align*}
Since $n(L,G) = 1$ for any subgroup $L$, and the intermediate terms vanish modulo 4 by the inductive hypothesis, this collapses to $\operatorname{coeff}^L ( \omega([a,b]) ) |W(L)| \equiv \omega^L([a,b]) - \omega^G([a,b]) \pmod 4$. By \eqref{eq:parity_lock_1}, this difference is exactly $0 \pmod 4$. Thus, $\operatorname{coeff}^L ( \omega([a,b]) ) |W(L)| \equiv 0 \pmod 4$ for all intermediate subgroups $H < L < G$.
Finally, we evaluate the relation at $H$:
\begin{align*}
\operatorname{coeff}^H ( \omega([a,b]) ) |W(H)| \equiv \omega^H([a,b]) - \sum_{H < L} \operatorname{coeff}^L ( \omega([a,b]) ) |W(L)| n(H,L) \pmod 4.    
\end{align*}
Since all intermediate terms $\operatorname{coeff}^L ( \omega([a,b]) ) |W(L)|$ vanish modulo 4, and $n(H,G) = 1$, we obtain $\operatorname{coeff}^H ( \omega([a,b]) ) |W(H)| \equiv \omega^H([a,b]) - \omega^G([a,b]) \pmod 4$. Since both values strictly belong to $\{0, 2\} \pmod 4$ and are unequal by \eqref{eq:parity_lock_2}, their difference must be exactly $2 \pmod 4$. Since $\operatorname{coeff}^H ( \omega([a,b]) ) |W(H)| \equiv 2 \pmod 4$, it follows that $\operatorname{coeff}^H ( \omega([a,b]) ) |W(H)| \neq 0$. In this way we have proven that the coefficient $\operatorname{coeff}^H(\omega([a,b]))$ must be non-zero.
\end{proof}

\end{document}
